\section{The decomposition theorem}\label{sec:8.6}

When dark matter kicks a silicon nucleus, the electrons follow most of the time, and the recoil energy ends up in the lattice as phonons. In a small fraction of events (\cref{sec:6.1}) the electrons are left behind, and the same kick leaves an electron--hole pair. At \SI{50}{keV} the nucleus does not recoil freely: it leaves the lattice untouched in \SI{21}{\percent} of the events and creates two or more phonons in \SI{47}{\percent} (\cref{we:5-dw}). The electron density offers the same variety. A slow particle that couples to the electrons excites a pair with no phonon (\cref{chap:4}), a pair with a phonon below the direct gap (\cref{sec:6.5}), or no pair, since the electrons that follow a vibrating nucleus scatter as well (\cref{sec:5.5}). The literature computes each of these final states with its own formula and its own code. The sum of these separate rates is the total rate only if no final state is counted twice and none is left out. \Cref{we:2-twophoton} and \cref{sec:6.5} have already met a double count: above the direct gap, the phonon-assisted absorption of light contains a cascade that direct absorption counts once more.

To make the sum count each final state once we split the final states of every response into three sectors. The derivations are in \cref{app:D.28}.

\subsection{Three sectors of final states}

The Born--Oppenheimer approximation of \cref{sec:5.1} already suggests the first cut. Every state of the form $\Psi^{\mathrm{nuc}}(u)\,\Psi_0(r;u)$ lies in the range of the projector
\begin{equation}
  \mathcal{P}_{\mathrm{BO}} = \int\dd u\;|\Psi_0(u)\rangle\langle\Psi_0(u)|\otimes|u\rangle\langle u| ,
  \label{eq:8.6-PBO}
\end{equation}
where $|\Psi_0(u)\rangle$ acts on the electrons and $|u\rangle$ on the nuclei. In a state of this form the electrons are in the adiabatic ground state at each configuration $u$ of the nuclei, and $\Psi^{\mathrm{nuc}}(u)$ is any nuclear wavefunction. Phonons, multiphonon states and a freely recoiling nucleus are all in this subspace; electron--hole pairs are not. The subspace is not exactly invariant under the Hamiltonian, however. The nuclear kinetic energy differentiates $\Psi_0(r;u)$ with respect to $u$, and the nonadiabatic couplings \cref{eq:5.1-N} carry part of the state into the excited adiabatic states. For zero-point motion the leak out of the subspace is the derivative coupling $\alpha^2m_e\,\kappa_{\mathrm{BO}}^3$ of the Born--Oppenheimer hierarchy (\cref{sec:6.1}).

A final state must have a definite sector, and a projector that leaks does not give it one. Panati, Spohn and Teufel showed how to repair the leak order by order~\cite[Thm.~2.1]{panati2003space}. Their small parameter is $(m_e/M)^{1/2} = \kappa_{\mathrm{BO}}^2$, the ratio of a nuclear to an electronic velocity at equal kinetic energy. A unitary $U_n$, built as a power series in it, rotates $\mathcal{P}_{\mathrm{BO}}$ into a \emph{superadiabatic} projector
\begin{equation}
  \mathcal{P}_{\mathrm{ad}} = U_n\,\mathcal{P}_{\mathrm{BO}}\,U_n^\dagger , \qquad
  \bigl\|(1 - \mathcal{P}_{\mathrm{ad}})\,H\,\mathcal{P}_{\mathrm{ad}}\bigr\| = O\bigl(\kappa_{\mathrm{BO}}^{2(n+1)}\bigr)\,\alpha^2m_e ,
  \label{eq:8.6-leak}
\end{equation}
on states below any fixed energy~\cite[Sect.~4.5]{panati2003space}. The order $n = 0$ is \cref{eq:8.6-PBO} itself. Each further order removes one more power of $\kappa_{\mathrm{BO}}^2$. The first correction tilts the electronic state in the direction of the nuclear velocity, $\Psi_0 \to \Psi_0 + \sum_{j\neq0}\Psi_j\,\ii\,\vect{d}_{j0}\cdot\vect{p}/[M_\kappa(E_j - E_0)]$, with $\vect{d}_{j0}$ the derivative-coupling vector \cref{eq:6.1-d} and $\vect{p}$ the momentum of the nucleus: the electron cloud of a moving nucleus lags behind it and is slightly polarized. \Cref{app:D.28} gives the construction and its hypotheses.

The range of $\mathcal{P}_{\mathrm{ad}}$ is the \emph{electron-following sector}. Its complement $1 - \mathcal{P}_{\mathrm{ad}}$ contains every state with an electronic excitation, and it splits once more by what happens to the lattice. Measured from the initial state, the number of phonons can stay the same or change. We write
\begin{equation}
  1 = \mathcal{P}_{\mathrm{ad}} + \mathcal{P}_0 + \mathcal{P}_+ ,
  \label{eq:8.6-partition}
\end{equation}
with $\mathcal{P}_0$ the \emph{zero-phonon electronic sector}, an electronic excitation with no net change of the phonon number, and $\mathcal{P}_+$ the \emph{phonon-dressed electronic sector}, an electronic excitation with a nonzero net change of the phonon number. At $T = 0$ phonons can only be emitted, and the change is at least one phonon created; a warm target can also give up phonons, as at the absorption thresholds of \cref{sec:6.5}. The phonon number here is the phonon number of the harmonic reference lattice of the excited electronic state. The excited surfaces are displaced, and a zero-phonon state of an excited surface is a coherent superposition of phonons of the ground surface, as the displaced oscillator of \cref{sec:5.6} showed. At $T = 0$, $\mathcal{P}_0$ holds the zero-phonon line of every electronic transition and $\mathcal{P}_+$ its sideband. The three projectors are orthogonal and sum to one, so every final state has one label. The lattice regime $\mathcal{K}$ of \cref{sec:8.5} labels the phonon content within a sector: one phonon, several or a free recoil inside $\mathcal{P}_{\mathrm{ad}}$, and the phonons that dress the excitation inside $\mathcal{P}_+$. The sector and the lattice regime are nested labels.

The response of each sector follows from \cref{eq:4.2-Sab} by inserting the projector on both sides of the energy delta function,
\begin{equation}
  S^A_{ab}(\vect{q},\omega) = \frac{2\pi}{\Vol}\sum_ip_i\,\bigl\langle i\big|O_a(\vect{q})^\dagger\,\mathcal{P}_A\,\delta(\omega + E_i - H)\,\mathcal{P}_A\,O_b(\vect{q})\big|i\bigr\rangle , \qquad A \in \{\mathrm{ad}, 0, +\} ,
  \label{eq:8.6-sector}
\end{equation}
with the initial states in the electron-following sector, as they are at $T \ll \Eg$. If $\mathcal{P}_A$ commutes with $H$, the delta function can be written in eigenstates that each lie in one sector, $S^A_{ab}$ is the sum over the final states of sector $A$, and $\sum_AS^A_{ab} = S_{ab}$ exactly. Each $S^A_{ab}$ is then a positive semidefinite matrix by the argument of \cref{sec:4.2}, because it is again a sum of outer products with positive weights. If $\mathcal{P}_A$ commutes with $H$ approximately, the cross terms $\mathcal{P}_A\,\delta(\omega + E_i - H)\,\mathcal{P}_B$ with $A \neq B$ do not vanish. They are small on energy bins wider than the leakage norm, and not on narrower ones. For a bin of width $\Delta\omega_{\mathrm{bin}}$ they are bounded by the leakage norm of \cref{eq:8.6-leak} divided by $\Delta\omega_{\mathrm{bin}}^2$, times the norms of $O_a|i\rangle$ and $O_b|i\rangle$ (\cref{eq:D.28-lemma}).
\begin{keyresult}[Decomposition theorem]
For a target with a gap, with the nuclear charges smeared over the finite size of the nuclei, at temperatures $T \ll \Eg$, and with the electronic ground state nondegenerate at every nuclear configuration that the vibrational state visits,
\begin{equation}
  S_{ab}(\vect{q},\omega) = S^{\mathrm{ad}}_{ab}(\vect{q},\omega) + S^0_{ab}(\vect{q},\omega) + S^+_{ab}(\vect{q},\omega) + \delta S_{ab} , \qquad
  \frac{|\delta S_{ab}|}{S^{\mathrm{bin}}} = O\Bigl(\kappa_{\mathrm{BO}}^{2(n+1)}\,\frac{\alpha^2m_e}{\Delta\omega_{\mathrm{bin}}}\Bigr) ,
  \label{eq:8.6-decomp}
\end{equation}
for every pair of target operators, with $S^{\mathrm{bin}}$ the response averaged over an energy bin of width $\Delta\omega_{\mathrm{bin}}$ and $n$ the order of the superadiabatic projector. The electron-following sector and the sum of the two electronic sectors are positive semidefinite matrices. The split between $S^0$ and $S^+$ is unique to relative order $\theta_{\mathrm{ep}}$ inside the off-shell window, where each is positive; outside it the split is fixed by the Dyson convention of \cref{eq:8.6-dyson}, and near a resonance only the sum of the two is positive.
\end{keyresult}
For bins of \SI{1}{meV}, the scale of phonon spectra, the factor $\alpha^2m_e/\Delta\omega_{\mathrm{bin}}$ is $2.7\times10^4$, and the bound becomes useful from $n = 2$ on, where it is $2\times10^{-3}$. The finite size of the nuclei keeps the electronic Hamiltonian a smooth function of the nuclear positions. The adiabatic construction requires this smoothness, and point charges violate it at coincidence~\cite[Sect.~2]{spohn2001adiabatic}; the nuclear charge radius changes nothing else in this book. The gap excludes metals and level crossings, and the temperature condition keeps the initial state in the electron-following sector. The theorem is proved for a finite crystal. It is open whether the theorem extends to the infinite crystal with constants that do not grow with the number of nuclei; this problem is stated with the other open problems at the end of this section. The theorem does not depend on which operators $O_a$ appear: the final state alone decides the sector.

\subsection{Densities in each sector}

Inside the electron-following sector an electronic operator acts on the nuclei alone. The electrons are in $\Psi_0(r;u)$, and the electron density reaches the sector only through its expectation value at each configuration,
\begin{equation}
  \mathcal{P}_{\mathrm{ad}}\,n_e(\vect{q})\,\mathcal{P}_{\mathrm{ad}} \simeq \bar\rho_e(\vect{q};u) \equiv \bigl\langle\Psi_0(u)\big|n_e(\vect{q})\big|\Psi_0(u)\bigr\rangle ,
  \label{eq:8.6-follow}
\end{equation}
up to corrections of relative order $\kappa_{\mathrm{BO}}^2$ from the tilt of the state. The function $\bar\rho_e(\vect{q};u)$ is the Fourier component of the ground-state electron density of the crystal frozen in configuration $u$, an operator on nuclear wavefunctions like the nuclear density $n_\kappa(\vect{q})$ itself. At $u = 0$ it is the static density, nonzero only at reciprocal-lattice vectors: Bragg scattering. Its term linear in $u$ creates one phonon, and its higher terms several. A dark-matter particle that couples only to electrons therefore excites phonons, through the electrons that are carried along by the vibrating nuclei. If each nucleus dragged a rigid cloud with form factor $f_\kappa(\vect{q})$, we would have $\bar\rho_e = \sum_If_{\kappa(I)}(\vect{q})\,\ee^{\ii\vect{q}\cdot\vect{X}_I}$, and the electron channel would see the lattice with the weight $|f_\kappa|^2$ of the nuclear channel. In a real crystal the cloud deforms. At long wavelength and zero macroscopic field, the electrons that follow nucleus $\kappa$ carry the charge $Z_\kappa\mathbf{1} - Z^*_\kappa$, a tensor, so that nucleus and electrons together carry the Born effective charge of \cref{eq:5.5-Z}. In silicon $Z^*_\kappa = 0$ and all fourteen electrons follow; in GaAs a fraction is left behind, and the lattice is seen by a dark photon through $Z^*_\kappa$, screened by $\varepsilon_\infty$ (\cref{sec:8.4}).

The electronic sectors are reached through the transition densities between adiabatic states. A final state $\Psi^{\mathrm{nuc}\prime}\Psi_j$ on the excited surface $j$, with $\Psi^{\mathrm{nuc}\prime}$ one of its vibrational states, is reached from $\Psi^{\mathrm{nuc}}\Psi_0$ with the amplitude
\begin{equation}
  \begin{gathered}
  \bigl\langle\Psi^{\mathrm{nuc}\prime}\Psi_j\big|n_e(\vect{q})\big|\Psi^{\mathrm{nuc}}\Psi_0\bigr\rangle = \bigl\langle\Psi^{\mathrm{nuc}\prime}\big|\rho_{j0}(\vect{q};u)\big|\Psi^{\mathrm{nuc}}\bigr\rangle , \\
  \rho_{j0}(\vect{q};u) = \rho_{j0}(\vect{q};0) + \sum_{\vect{R}\kappa}\vect{u}_{\vect{R}\kappa}\cdot\nabla_{\vect{R}\kappa}\rho_{j0} + \cdots ,
  \end{gathered}
  \label{eq:8.6-electronic}
\end{equation}
where $\rho_{j0}(\vect{q};u) = \langle\Psi_j(u)|n_e(\vect{q})|\Psi_0(u)\rangle$ is the transition density of the frozen crystal. The constant term is the Condon approximation of \cref{sec:5.6}. It weights the vibrational states by the Franck--Condon overlaps $|\langle\Psi^{\mathrm{nuc}\prime}|\Psi^{\mathrm{nuc}}\rangle|^2$: the zero-phonon line, in $\mathcal{P}_0$, with the weight $\ee^{-S_{\mathrm{HR}}}$ of \cref{eq:5.6-FC}, and the replicas, in $\mathcal{P}_+$. For an electron--hole pair in a crystal the surfaces of the ground and excited states are displaced by almost nothing, since one pair spread over many cells exerts a negligible force on each nucleus, and the zero-phonon line carries nearly all of the Condon term. The linear term, the Herzberg--Teller term of molecular spectroscopy, creates a phonon even without a displacement. In a crystal it does more: at $u \neq 0$ the frozen lattice is not periodic, and $\rho_{j0}$ connects pairs of different crystal momentum that the static lattice keeps apart. The linear term is phonon-assisted electron excitation in adiabatic form, the term that the special-displacement supercell of \cref{sec:6.5} computes. Being adiabatic, the linear term misses the phonon energy in the final state, a correction of relative order $\omega_{\vect{Q}\nu}/\delta_{\mathrm{off}}$; the electron--phonon form \cref{eq:6.5-eps2PA} restores this energy.

The nuclear density cannot reach the electronic sectors in the Born--Oppenheimer approximation. Its amplitude between the same two states is $\int\dd u\,\Psi^{\mathrm{nuc}\prime*}(u)\,\ee^{\ii\vect{q}\cdot\vect{X}_I}\Psi^{\mathrm{nuc}}(u)\,\allowbreak\langle\Psi_j(u)|\Psi_0(u)\rangle$, and the electronic overlap vanishes at every $u$ because the adiabatic states are orthogonal. The amplitude comes from the tilt of the initial and final states by the superadiabatic correction. At its first order (\cref{app:D.28}),
\begin{equation}
  \bigl\langle\widetilde{\Psi^{\mathrm{nuc}\prime}\Psi_j}\big|\ee^{\ii\vect{q}\cdot\vect{X}_I}\big|\widetilde{\Psi^{\mathrm{nuc}}\Psi_0}\bigr\rangle = -\ii\,\frac{\vect{q}\cdot\vect{d}^{I}_{j0}}{M_\kappa(E_j - E_0)}\,\bigl\langle\Psi^{\mathrm{nuc}\prime}\big|\ee^{\ii\vect{q}\cdot\vect{X}_I}\big|\Psi^{\mathrm{nuc}}\bigr\rangle ,
  \label{eq:8.6-boost}
\end{equation}
with $\vect{d}^I_{j0}$ the derivative coupling with respect to the position of nucleus $I$ of species $\kappa$, and the tildes marking the tilted states. Two terms proportional to the gradient of the nuclear wavefunctions cancel between the initial and the final tilt, and only the momentum $\vect{q}$ of the kick survives. The electronic factor is the sudden amplitude $\vect{v}_N\cdot\vect{d}_{j0}/(\ii\Delta E_j)$ of \cref{eq:6.1-sudden1}, with $\vect{v}_N = \vect{q}/M_\kappa$: the electrons see a nucleus that suddenly moves. The nuclear factor is the kick of \cref{sec:5.6}, with its Debye--Waller and multiphonon structure. The Migdal effect is purely nonadiabatic. Its zero-phonon term lies in $\mathcal{P}_0$, and the terms in which the kick also creates phonons, or a free recoil, lie in $\mathcal{P}_+$. In the truncation label of \cref{sec:8.5}, which counts powers in the rate, a Migdal rate carries the derivative coupling squared and is at $n_B = 2$. The interference of a Migdal amplitude with an electron recoil to the same final state is linear in the derivative coupling, at $n_B = 1$.

The cross response $S_{e\kappa}$ is nonzero in every sector whenever the dark matter couples to electrons and nuclei at once, as a dark photon does, and the next subsection shows that it cannot be dropped as small without violating a sum rule.

In practice each sector comes from a different chain of codes, and the three calculations must use one set of phonon energies, vertices and band energies, or the sum of the three sector responses is the response of no single Hamiltonian.

\subsection{Partial sum rules}

The first-moment identity \cref{eq:4.3-firstmoment} holds sector by sector, because a projector that commutes with $H$ can be carried through the double commutator. Restricted to the electron-following sector, both the electron and the nuclear density are functions of the nuclear positions, and the only term of the Hamiltonian that fails to commute with them is the kinetic energy of the nuclei, exactly as for $n_\kappa$ in \cref{sec:4.3}. For the electron density, \cref{eq:8.6-follow} turns the double commutator into a sum over nuclei of squared gradients (\cref{app:D.28}),
\begin{equation}
  \mathscr{M}^1_{ee}\big|_{\mathrm{ad}}(\vect{q}) = \frac1\Vol\Bigl\langle\sum_{\vect{R}\kappa}\frac{|\nabla_{\vect{R}\kappa}\bar\rho_e(\vect{q};u)|^2}{M_\kappa}\Bigr\rangle\bigl[1 + O(\kappa_{\mathrm{BO}}^2)\bigr]
  \;\simeq\; q^2\sum_\kappa\frac{\nd_\kappa\,|f_\kappa(\vect{q})|^2}{M_\kappa} ,
  \label{eq:8.6-fsum-e}
\end{equation}
where the second form holds for rigid clouds. Measured against the f-sum $\nd_eq^2/m_e$, the electron-following first moment is of relative size $Z_\kappa m_e/M_\kappa$ at small $q$, where the whole cloud follows. The relative size is $2.73\times10^{-4}$ for silicon with all fourteen electrons, the number that \cref{sec:4.3} found for the nuclear share of the f-sum of the total charge, and for the same reason: either way, $Z_\kappa$ units of charge are attached to a mass $M_\kappa$. In powers of the small parameter it is $\kappa_{\mathrm{BO}}^4 = m_e/M$. The two electronic sectors together therefore carry the f-sum up to a part in $10^3$, and the f-sum test of a clamped-nuclei calculation (\cref{sec:4.3}) applies to them unchanged.

For the nuclear density the electronic sectors carry the Migdal events, and their share follows from \cref{eq:8.6-boost}. Summed over the vibrational states $\Psi^{\mathrm{nuc}\prime}$, which form a complete set, the probability that a kick of nucleus $I$ leaves the electrons excited is
\begin{equation}
  P_{\mathrm{Mig}}(\vect{q}) = \sum_{j\neq0}\frac{|\vect{q}\cdot\vect{d}^I_{j0}|^2}{M_\kappa^2(E_j - E_0)^2} ,
  \label{eq:8.6-PMig}
\end{equation}
at leading order, the sudden probability of \cref{sec:6.1} with $\vect{v}_N = \vect{q}/M_\kappa$. Each Migdal event deposits the electronic energy $E_j - E_0$ plus, on average, the recoil energy $E_R = q^2/2M_\kappa$ of \cref{eq:5.6-sum}. With $\vect{d}_{j0} = \vect{F}_{j0}/(E_0 - E_j)$ from \cref{eq:5.1-dcoup}, the electronic part of the deposit adds up to $\vect{q}\cdot\delta M_\kappa\cdot\vect{q}/2M_\kappa^2$, and the partial sum rules per nucleus read (\cref{app:D.28})
\begin{equation}
  \begin{aligned}
  \mathscr{M}^1_{\kappa\kappa}\big|_{0+} &= \nd_\kappa\Bigl[P_{\mathrm{Mig}}\,\frac{q^2}{M_\kappa} + \frac{\vect{q}\cdot\delta M_\kappa\cdot\vect{q}}{M_\kappa^2}\Bigr] , \qquad
  \mathscr{M}^1_{\kappa\kappa}\big|_{\mathrm{ad}} = \nd_\kappa\,(1 - P_{\mathrm{Mig}})\;\vect{q}\cdot(M_\kappa + \delta M_\kappa)^{-1}\cdot\vect{q} , \\
  \delta M_\kappa &= 2\sum_{j\neq0}\frac{\Real\bigl[\vect{F}_{0j}\,\vect{F}_{j0}\bigr]}{(E_j - E_0)^3} , \qquad \vect{F}_{j0} = \bigl\langle\Psi_j\big|\nabla_{\vect{R}\kappa}V_{eN}\big|\Psi_0\bigr\rangle ,
  \end{aligned}
  \label{eq:8.6-fsum-k}
\end{equation}
to first order in $P_{\mathrm{Mig}}$ and $\delta M_\kappa/M_\kappa$, with $0+$ marking the two electronic sectors together. The tensor $\delta M_\kappa$ is the \emph{electron drag mass}: a nucleus that moves slowly drags part of its electron cloud along and vibrates as if it were heavier. The solvable model of \cref{app:D.13} found this effect at order $m_e/M$. The same term appears in the second-order Born--Oppenheimer Hamiltonian of Panati, Spohn and Teufel~\cite[Eq.~(40), last term]{panati2003space}. For an isolated atom the sum is fixed by the Thomas--Reiche--Kuhn rule \cref{eq:4.3-TRK}, and $\delta M = N_em_e$: the atom recoils with its electrons, as it must (\cref{app:D.28}). In a crystal the f-sum rule fixes the sum of the drag masses over all nuclei, exactly $14\,m_e$ per silicon atom (\cref{eq:12.5-trace}), and this sum gives $\delta M_\kappa/M_\kappa \approx 2.7\times10^{-4}$. A kick on one nucleus probes instead the on-site drag mass, which no sum rule fixes: \cref{sec:12.5-trace} evaluates its valence part as $18.6\,m_e$, between 11.8 and $24.0\,m_e$ in the approximations behind that evaluation.

\Cref{eq:8.6-fsum-k} refines \cref{eq:5.6-sum}. That rule says that the mean energy given to the lattice is $E_R$ on any potential-energy surface, and it is a statement within the Born--Oppenheimer approximation. Beyond the Born--Oppenheimer approximation the total first moment is still $E_R$ per nucleus, by \cref{eq:4.3-species}, but the lattice receives less. The electrons take $P_{\mathrm{Mig}}E_R + \vect{q}\cdot\delta M_\kappa\cdot\vect{q}/2M_\kappa^2$, and the lattice the rest of $E_R$, which is the recoil of a nucleus of mass $M_\kappa + \delta M_\kappa$ in the events that leave the electrons alone. In silicon the two electronic terms stay below $10^{-3}$ of $E_R$ for recoil energies up to about \SI{16}{eV}, a range that covers the phonon and multiphonon regime. Corrections below $10^{-3}$ are far below the accuracy of any multiphonon calculation; at \SI{500}{eV} the weight $1 - P_{\mathrm{Mig}}$ is a correction of \SI{3}{\percent}. The bookkeeping, however, is exact, and \cref{sec:12.5} carries it through.

The total first moment of the electron--nucleus interference vanishes, \cref{eq:4.3-species}, but the first moments of the individual sectors do not. In the electron-following sector the same steps that led to \cref{eq:8.6-fsum-e} give
\begin{equation}
  \mathscr{M}^1_{e\kappa}\big|_{\mathrm{ad}}(\vect{q}) = \frac{1}{\Vol}\Bigl\langle\sum_{I\in\kappa}\frac{\nabla_I\bar\rho_e(\vect{q};u)^\dagger\cdot\ii\vect{q}\,\ee^{\ii\vect{q}\cdot\vect{X}_I}}{M_\kappa}\Bigr\rangle \simeq \frac{q^2\,\nd_\kappa\,f_\kappa(\vect{q})^*}{M_\kappa} , \qquad
  \mathscr{M}^1_{e\kappa}\big|_{0+} = -\,\mathscr{M}^1_{e\kappa}\big|_{\mathrm{ad}} ,
  \label{eq:8.6-fsum-ek}
\end{equation}
to leading order in $\kappa_{\mathrm{BO}}$. Electrons that follow the nucleus move in phase with it and carry a positive energy-weighted interference. The electronic sectors carry the same weight with the opposite sign. The opposite signs of the two parts are the sign change that \cref{sec:4.3} required of the interference spectrum, now located: the electron-following and the electronic parts of $S_{e\kappa}$ are mirror images in their first moment. If a model keeps the Born-charge interference and drops the electronic parts of $S_{e\kappa}$, it violates the sum rule by the full Born-charge term.

These rules give prescriptions for adding rates. If a multiphonon $S_{\kappa\kappa}$ normalized to the full nuclear f-sum is added to a Migdal rate, the electronic share is counted twice; the electron-following part must be depleted by $P_{\mathrm{Mig}}$ and computed with the mass $M_\kappa + \delta M_\kappa$. A resonant sideband is never added to a zero-phonon line that already contains its cascade (\cref{sec:2.4,sec:6.5}); the next subsection makes this rule precise. A measured loss function or optical conductivity contains both electronic sectors, and a computed sideband is never added to it. And when the dark matter couples to both species, the electron and nuclear amplitudes to the same final state interfere in every sector, with the interference bounded by the positivity of \cref{eq:4.2-positivity}.

\subsection{Zero-phonon line and phonon sideband}

The superadiabatic projector commutes with $H$ to the order of \cref{eq:8.6-leak} because a gap separates the electronic states. Nothing comparable, however, separates $\mathcal{P}_0$ from $\mathcal{P}_+$. A zero-phonon pair at energy $E$ and a one-phonon state at the same energy are both electronic excitations, and the electron--phonon coupling $H_{\mathrm{ep}}$ of \cref{sec:6.2} turns one into the other. Whether the split between $\mathcal{P}_0$ and $\mathcal{P}_+$ means anything depends on where the energy deposited $\omega$ lies relative to the zero-phonon states at the same momentum transfer.

Below the zero-phonon states the answer is simple. Take $\omega$ inside the off-shell window, below the bottom $E_{\mathrm{min}}(q)$ of the zero-phonon pair continuum by a margin $\delta_{\mathrm{w}}$ large compared with the phonon energies, the widths and $|g|$ (the window is made precise in \cref{sec:13.6}). No zero-phonon state exists at that energy, and every final state contains a phonon. An exact final state $|\tilde f\rangle$ is a one-phonon state $|f\rangle$ dressed by a virtual admixture of zero-phonon pairs $m$, each off shell by at least $\delta_{\mathrm{w}}$,
\begin{equation*}
  |\tilde f\rangle = |f\rangle + \sum_m|m\rangle\,\frac{\langle m|H_{\mathrm{ep}}|f\rangle}{E_f - E_m} + \cdots , \qquad
  \bigl\|\mathcal{P}_0^{(0)}|\tilde f\rangle\bigr\|^2 = O(\theta_{\mathrm{ep}}) ,
\end{equation*}
with $\mathcal{P}_0^{(0)}$ the zero-phonon projector of the uncoupled reference and $\theta_{\mathrm{ep}}$ the off-shell parameter \cref{eq:8.5-theta}. The perturbation series converges, since no denominator comes close to zero, and the state can be followed continuously from $g = 0$. Every reasonable definition of its label agrees: the phonon number of its leading component, the number obtained after a Schrieffer--Wolff transformation that removes $H_{\mathrm{ep}}$ to a given order, and the number of phonons that eventually leave. The three definitions differ by the weight $O(\theta_{\mathrm{ep}})$ of the virtual admixture, and the whole response at $\omega$ belongs to $\mathcal{P}_+$ to that accuracy (\cref{app:D.28}). In silicon, at the absorption edge, $\theta_{\mathrm{ep}} = 0.010$ to 0.016 (\cref{sec:8.5}). The bare projector would not do. Inserted in \cref{eq:8.6-sector}, $\mathcal{P}_0^{(0)}$ would pick out the virtual admixture through which the whole phonon-assisted amplitude passes, and assign to the zero-phonon sector a weight as large as the rate itself. Sector projectors must commute with $H$, and here commuting with $H$ requires dressed projectors.

Above $E_{\mathrm{min}}(q)$ the zero-phonon pair is real, and the label becomes a matter of time. A pair $m$ created at energy $\omega$ lies inside the continuum of one-phonon states $f$ and decays into it by emitting a phonon. Its amplitude follows from the level-and-band model of \cref{we:4-level}, with the zero-phonon pair as the level and the one-phonon states as the band. Let $\rho_m = \langle m|O_b|i\rangle$ be the zero-phonon vertex, $g_{fm} = \langle f|H_{\mathrm{ep}}|m\rangle$ the electron--phonon matrix element built from the $g_{mn\nu}$ of \cref{eq:6.2-g}, and $g_{fm}\mathcal{A}^{\mathrm{off}}_f$ the amplitude with which the probe reaches $f$ by every path that does not pass through $m$ on shell: the phonon emitted first, hole orderings, and nonresonant pairs. The amplitude into $f$ at $\omega = E_f$ is then exactly (\cref{app:D.28})
\begin{equation}
  A_f(\omega) = g_{fm}\Bigl[\mathcal{A}^{\mathrm{off}}_f + \GF_m(\omega)\bigl(\rho_m + \Upsilon_m(\omega)\bigr)\Bigr] , \qquad
  \GF_m(\omega) = \frac{1}{\omega - E_m - \Sigma_m(\omega)} ,
  \label{eq:8.6-amp}
\end{equation}
with the self-energy of the pair $\Sigma_m(\omega) = \sum_f|g_{fm}|^2/(\omega - E_f + \ii0^+)$ and the vertex correction $\Upsilon_m(\omega) = \sum_f|g_{fm}|^2\mathcal{A}^{\mathrm{off}}_f/(\omega - E_f + \ii0^+)$. The width of the pair is $\varGamma_m(\omega) = -2\Imag\Sigma_m(\omega)$, its part due to phonon branch $\nu$ is $\varGamma^{(\nu)}_m$, and the quasiparticle energy $\tilde E_m$ and residue $z_{\mathrm{qp}}(m)$ follow from $\Sigma_m$ as in \cref{eq:4.7-qp}. The response is $(2\pi/\Vol)\sum_f|A_f|^2\delta(\omega - E_f)$. Near $\tilde E_m$ the squared propagator $|\GF_m|^2$ is a Lorentzian of width $\varGamma_m$, and \cref{eq:2.4-split} splits it into a delta function and a finite part. We adopt the convention that assigns the delta function to the zero-phonon sector and everything else to the phonon-dressed sector.
\begin{keyresult}[Dyson convention]
For $\varGamma_m$ small compared with the detuning $\delta_{\mathrm{F}}$, the distance from $\tilde E_m$ to the nearest structure of the one-phonon continuum,
\begin{equation}
  \begin{aligned}
  S^0(\vect{q},\omega) &= \frac{2\pi}{\Vol}\sum_mz_{\mathrm{qp}}(m)\,|\tilde\rho_m|^2\,\delta(\omega - \tilde E_m) , \\
  S^+(\vect{q},\omega) &= \frac{2\pi}{\Vol}\sum_m\sum_f|g_{fm}|^2\,\delta(\omega - E_f)\Bigl[|\rho_m|^2\operatorname{fp}\frac{1}{(\omega - \tilde E_m)^2} \\
  &\qquad + 2\Real(\mathcal{A}^{\mathrm{off}*}_f\rho_m)\operatorname{PV}\frac{1}{\omega - \tilde E_m} + |\mathcal{A}^{\mathrm{off}}_f|^2\Bigr] ,
  \end{aligned}
  \label{eq:8.6-dyson}
\end{equation}
where $|\tilde\rho_m|^2 = |\rho_m|^2 + 2\Real[\rho_m^*\operatorname{PV}\Upsilon_m(\tilde E_m)]$ contains the virtual-phonon correction of the vertex. Inside the off-shell window no $\tilde E_m$ lies at $\omega$, the delta term is absent, and $S^+$ is the ordinary second-order rate.
\end{keyresult}
The delta part is a real zero-phonon pair, which is handed to the cascade of \cref{chap:18} as a real pair; the cascade lets the pair emit its first phonon (\cref{sec:8.7}). The finite-part and principal-value terms are the sideband, a pair and a phonon created together. The finite part is negative against a narrow positive test function, as \cref{sec:2.4} found, so $S^+$ alone is not a positive spectrum near $\tilde E_m$.

The convention preserves the sum of the sectors. Against any test function smooth on the scale $\delta_{\mathrm{F}}$, such as a detector resolution or the integral over $\vect{q}$ and the velocity distribution, $S^0 + S^+$ reproduces the exact response up to terms of order $g^4$ and $\varGamma_m^2/\delta_{\mathrm{F}}^2$ (\cref{app:D.28}). The step that makes the weights add up is a per-pair zeroth-moment identity. Differentiating the Kramers--Kronig form of $\Real\Sigma_m$ with \cref{eq:C.4-dPV} gives $\partial_\omega\Real\Sigma_m = -(2\pi)^{-1}\operatorname{fp}\!\int\dd\omega'\,\varGamma_m(\omega')/(\omega' - \omega)^2$, and since $z_{\mathrm{qp}} = 1 + \partial_\omega\Real\Sigma_m + O(g^4)$,
\begin{equation}
  z_{\mathrm{qp}}(m) + \frac{1}{2\pi}\operatorname{fp}\!\int\dd\omega'\,\frac{\varGamma_m(\omega')}{(\omega' - \tilde E_m)^2} = 1 + O(g^4) .
  \label{eq:8.6-zsum}
\end{equation}
The weight that the zero-phonon line loses to its residue is exactly the finite-part weight of the sideband, the first term of $S^+$ integrated over $\omega$. The level of \cref{we:4-level} shows that the residue is not a probability. There $z_{\mathrm{qp}} = 1.0064$, and the finite part of the flat band is $-2\varGamma/(\pi E_{\mathrm{w}}) = -0.0064$: the line carries more than the full weight and the off-shell part returns the excess. Inside a continuum whose width grows with energy, the residue also exceeds one, and a negative finite-part weight in a computed sideband is no error. The identity concerns the self-energy only; the vertex correction in $|\tilde\rho_m|^2$ redistributes weight that no zeroth-moment rule fixes, and there is no global zeroth-moment rule for $n_e$.

The split itself is a convention, and its ambiguity has a size. Three choices of convention are natural. The first is the power of the residue, which a time-domain definition by the surviving amplitude raises from $z_{\mathrm{qp}}$ to $z_{\mathrm{qp}}^2$. The second is the shape of the line, a Lorentzian or any normalized kernel no wider than $\varGamma_m$ in place of the delta function. The third is a window $|\omega - \tilde E_m| < x_{\mathrm{w}}$, with $\varGamma_m \ll x_{\mathrm{w}} \le \delta_{\mathrm{F}}$, all of whose weight is assigned to the line. If two conventions differ from \cref{eq:8.6-dyson} only in these ways, the zero-phonon weights $w_0$ that they assign in a bin around $\tilde E_m$ differ by (\cref{app:D.28})
\begin{equation}
  \frac{|\Delta w_0|}{w_0} \le \frac{5}{2\pi}\,\frac{\hat\varGamma_m}{\delta_{\mathrm{F}}} + \frac{\hat\varGamma_m}{\pi x_{\mathrm{w}}} + \Bigl(1 + \frac{|\mathcal{A}^{\mathrm{off}}|\,\delta_{\mathrm{F}}}{|\rho_m|}\Bigr)^2\frac{\hat\varGamma_mx_{\mathrm{w}}}{\delta_{\mathrm{F}}^2} ,
  \label{eq:8.6-ambiguity}
\end{equation}
with $\hat\varGamma_m$ the largest width within $\delta_{\mathrm{F}}$ of the line. At $x_{\mathrm{w}} \approx \delta_{\mathrm{F}}$ the first two terms add up to $1.1\,\hat\varGamma_m/\delta_{\mathrm{F}}$, and the bound is of order $\hat\varGamma_m/\delta_{\mathrm{F}}$ times a number of order one when the off-shell amplitude is no larger than $|\rho_m|/\delta_{\mathrm{F}}$. The split is well defined where the pair lives long compared with $1/\delta_{\mathrm{F}}$ and becomes meaningless as $\varGamma_m$ approaches $\delta_{\mathrm{F}}$. The bound does not extend to every convention. If the split is defined by the number of phonons present after an infinite time, every pair inside the continuum eventually decays, the zero-phonon sector is empty, and the split moves by an amount of order one. The sector label belongs to a time window between $1/\delta_{\mathrm{F}}$ and $1/\varGamma_m$, and the Dyson convention is the choice that makes the totals and the cascade consistent. 

Two mixed conventions are common and wrong. Adding a Lorentzian zero-phonon line to the finite-part sideband counts the off-shell weight $(\varGamma_m/2\pi)|\rho_m|^2\operatorname{fp}(1/x^2)$ twice, an error of order $g^2$. Adding a Lorentzian line to the resummed sideband $|x + \ii\varGamma_m/2|^{-2}$ counts the delta function twice, an error of order one; this second error is the double counting of the cascade in \cref{sec:2.4}. Within a few widths of $E_{\mathrm{min}}(q)$, or of a singularity of the joint density of states, no test function is smooth on the scale $\delta_{\mathrm{F}}$, and only totals computed from the unexpanded amplitude \cref{eq:8.6-amp} are quoted.

There is one more term, and it changes no total. When the pair decays, its first decay product can interfere with the same product made directly, and this interference shifts the branching of the first decay among the channels. The delta function receives, besides $z_{\mathrm{qp}}|\tilde\rho_m|^2$, a contribution from the interference between the off-shell path $\mathcal{A}^{\mathrm{off}}_f$ and the decay of the real pair. Summed over the final states of one decay channel $\nu$, the zero-phonon events that decay first into channel $\nu$ occur with the probability (\cref{app:D.28})
\begin{equation}
  \frac{\varGamma^{(\nu)}_m}{\varGamma^{\mathrm{tot}}_m}\bigl(1 + \bar c^{\mathrm{F}} - c^{\mathrm{F}}_\nu\bigr) , \qquad
  c^{\mathrm{F}}_\nu = \frac{\varGamma^{\mathrm{tot}}_m\Imag(\rho_m^*\mathcal{A}^{\mathrm{off}}_\nu)}{|\rho_m|^2} , \qquad
  \bar c^{\mathrm{F}} = \sum_\nu\frac{\varGamma^{(\nu)}_m}{\varGamma^{\mathrm{tot}}_m}\,c^{\mathrm{F}}_\nu ,
  \label{eq:8.6-fano}
\end{equation}
with $\nu$ running over every decay channel of the pair, each phonon branch and the others, such as impact ionization, $\mathcal{A}^{\mathrm{off}}_\nu$ the off-shell amplitude on the shell of channel $\nu$, and $\varGamma^{\mathrm{tot}}_m = \sum_\nu\varGamma^{(\nu)}_m$ the total width. The vertex correction $\Upsilon_m$ then sums over all channels as well. The deviations $\bar c^{\mathrm{F}} - c^{\mathrm{F}}_\nu$ of the \emph{Fano coefficients} from their mean have a weighted average of zero, with the weights $\varGamma^{(\nu)}_m$, because $\nu$ runs over all channels. The deviations therefore change no total and no number of zero-phonon events; the coefficients themselves need not average to zero. The deviations change the branching of the first decay, at relative order $\varGamma_m|\mathcal{A}^{\mathrm{off}}|/|\rho_m| \sim \varGamma_m/\delta_{\mathrm{F}}$, which is the order of the cascade itself. The Fano term is the interference of a discrete state with a continuum that Fano analyzed for atomic autoionization~\cite{fano1961effects}. It is unrelated to the Fano factor of the number of electron--hole pairs (\cref{sec:18.2}), which measures the variance of a count. An incoherent cascade cannot generate the term, since it is a coherence between production and decay, and the response calculation must supply it.

In many crystals the Fano term vanishes. Let the target be invariant under inversion combined with time reversal, and let its number of electrons be even. The antiunitary operator $\mathsf{P}\mathsf{T}$ then squares to one, and every state of the electrons and phonons can be chosen invariant under it, with the phases of the phonon polarization vectors fixed to match. The density $n_e(\vect{q})$ and $H_{\mathrm{ep}}$ are invariant too, so every matrix element and every off-shell amplitude is real, and $\Imag(\rho_m^*\mathcal{A}^{\mathrm{off}}_\nu) = 0$ (\cref{app:D.28}). Silicon and germanium satisfy the hypothesis, with or without spin--orbit coupling, at the level of one-body self-energies. In silicon and germanium the level-shift matrix of a one-body self-energy is a multiple of the identity on each degenerate zero-phonon manifold, a Kramers pair or an irreducible representation, so the invariant basis also diagonalizes it. Level shifts from the electron--hole vertex can split such a manifold in a basis that is not invariant; then only the sum of the Fano terms over the branches vanishes, not each one. GaAs has no center of inversion, and its Fano coefficients are generically nonzero.

\subsection{Targets without a gap}

Aluminum has electron--hole pairs at every energy down to zero, and its phonons, up to \SI{40}{meV} (\cref{sec:6.4}), lie inside that continuum. The adiabatic construction of \cref{eq:8.6-leak} needs a gap between the electronic ground state and the excited states at every nuclear configuration, and a metal has none: as \cref{sec:5.1} found, the expansion in adiabatic states fails at the Fermi surface. A superconductor has a gap, but $2\Delta_{\mathrm{BCS}} = \SI{0.36}{meV}$ in aluminum lies far below the phonon energies, and the adiabatic hierarchy is reversed: the electrons are the slow degrees of freedom near the Fermi surface. A Dirac semimetal is gapless at its nodes. Magnetic insulators are a fourth case. Their charge excitations are gapped, but at clamped nuclei the magnons are electronic excitations whose gaps, set by the magnetic anisotropy, lie below most phonon energies, so that no gap separates them from the lattice. The construction then goes through with the whole low-energy spin manifold in place of the electronic ground state, so that magnons and phonons share the electron-following sector (the magnetic targets of \cref{chap:14}).

For the gapless targets the sectors must count quasiparticles instead. The low-energy excitations of a metal are the quasiparticles of \cref{sec:4.7} and the phonons, and Migdal's theorem (\cref{sec:6.3}) separates them with the small parameter $\omega_{\mathrm{D}}/E_{\mathrm{F}} \sim \kappa_{\mathrm{BO}}^2$. Suppose a dressing unitary $U_{\mathrm{qp}}$, of the kind that Landau's theory and the Migdal--Eliashberg equations construct order by order, maps bare electrons and phonons onto quasiparticles and dressed phonons with number operators $\op{N}_{\mathrm{qp}}$ and $\op{N}_{\mathrm{ph}}$. Let $\mathcal{P}^{\mathrm{qp}}_{n,m}$ be the joint spectral projectors of $U_{\mathrm{qp}}\op{N}_{\mathrm{qp}}U_{\mathrm{qp}}^\dagger$ and $U_{\mathrm{qp}}\op{N}_{\mathrm{ph}}U_{\mathrm{qp}}^\dagger$, with $n$ and $m$ the net numbers of quasiparticles and phonons created. These projectors are orthogonal and sum to one. They do not commute with $H$, since quasiparticles have a width $\varGamma_{\mathrm{qp}}$ and phonons decay into pairs. Split $H = H_{\mathrm{d}} + H_{\mathrm{off}}$, with $H_{\mathrm{d}} = \sum_{n,m}\mathcal{P}^{\mathrm{qp}}_{n,m}H\mathcal{P}^{\mathrm{qp}}_{n,m}$, and define the sector responses with $H_{\mathrm{d}}$, which commutes with every $\mathcal{P}^{\mathrm{qp}}_{n,m}$. Each is then positive, and they sum exactly to the response of $H_{\mathrm{d}}$. On a Gaussian bin of standard deviation $\Delta\omega_{\mathrm{bin}}$, the response of $H$ from one initial state $|i\rangle$, without the factor $2\pi/\Vol$, differs from the response of $H_{\mathrm{d}}$ by at most (\cref{app:D.28})
\begin{equation}
  |C^{\mathrm{cross}}_{ab}| \le \frac{\mathfrak{h}_{\mathrm{off}}}{\pi\,\Delta\omega_{\mathrm{bin}}^2}\,\bigl\|O_a|i\rangle\bigr\|\,\bigl\|O_b|i\rangle\bigr\| , \qquad
  \mathfrak{h}_{\mathrm{off}} = \sup_{t'}\frac{\bigl\|H_{\mathrm{off}}\,\ee^{-\ii H_{\mathrm{d}}t'}\,O_b|i\rangle\bigr\|}{\bigl\|O_b|i\rangle\bigr\|} .
  \label{eq:8.6-bin}
\end{equation}
With the dressing in place, the bound needs only the few-quasiparticle state that the probe creates, not the whole crystal, and it holds for any orthogonal sectors; with the leakage of \cref{eq:8.6-leak} in place of $\mathfrak{h}_{\mathrm{off}}$ it is the remainder of the decomposition theorem. Relative to a typical binned value, $\|O|i\rangle\|^2/\Delta\omega_{\mathrm{bin}}$, the cross terms are of order $\mathfrak{h}_{\mathrm{off}}/\Delta\omega_{\mathrm{bin}}$. For a state that decays by the golden rule into a continuum of width $\Delta E_{\mathrm{c}}$, $\mathfrak{h}_{\mathrm{off}}^2 \approx \varGamma\,\Delta E_{\mathrm{c}}/2\pi$, which is much larger than $\varGamma^2$. The bound therefore guarantees small cross terms only on bins wide compared with $(\varGamma\,\Delta E_{\mathrm{c}})^{1/2}$. Without any dressing, $U_{\mathrm{qp}} = 1$, the bound fails: $H_{\mathrm{off}}$ acts on the filled Fermi sea as well as on the quasiparticles the probe creates, and this action on the Fermi sea, the disconnected part, makes $\mathfrak{h}_{\mathrm{off}}$ grow with the size of the crystal (\cref{app:D.28}). A bound has to control the connected part, in which $H_{\mathrm{off}}$ acts on the probe's excitations. Acting on a quasiparticle, $H_{\mathrm{off}}$ has two parts: the electron--phonon coupling, which adds or removes a phonon, and the part of the Coulomb interaction that turns the quasiparticle into two particles and a hole at fixed phonon number. The final states of the two parts differ in phonon number, so their squared norms add. The square of the connected part is then at most the largest eigenvalue of a band matrix, formed by the zeroth moments of the Fan--Migdal spectral function and by the analogous electron--electron sums over final states with two particles and a hole (\cref{app:D.28}). The electron--phonon entries are the sums $N^{-1}\sum|g_{mn\nu}|^2(2n_{\vect{Q}\nu} + 1)$ of \cref{sec:6.3} without their energy delta function. In aluminum, with $|g_{mn\nu}|^2$ taken constant over the band, the electron--phonon matrix alone has the eigenvalue $\int\dd\omega\,\alpha^2F/\mathcal{N}_{\mathrm{F}}\Omega = \SI{5.39}{meV}/\SI{0.193}{eV^{-1}} \approx \SI{0.03}{eV^2}$ (\cref{sec:6.4}), so the electron--phonon contribution to the connected part is at most about \SI{0.17}{eV}. The electron--electron entries cannot lower this eigenvalue, and they are not estimated.

The sectors of a superconductor follow the counting of \cref{sec:7.2}. A density or spin operator acting on the BCS state creates Bogoliubov quasiparticles in pairs, so $n \in \{0, 2\}$ at leading order. Phonons with no broken pair, $(0, m)$, are the analogue of the electron-following sector, and they are stable only below $2\Delta_{\mathrm{BCS}}$, since a phonon above it breaks a pair. Pair breaking, $(2, 0)$, computed with the conserving response of \cref{sec:7.2}, is the analogue of the zero-phonon sector, and phonon-assisted pair breaking, $(2, m \ge 1)$, the analogue of the phonon-dressed sector. The collective modes leave no sub-gap line in the density response: the phase mode is pushed to the plasma frequency, and the amplitude mode is not excited by the density in a particle--hole-symmetric metal (\cref{sec:7.2}). Assigning these collective modes, which are composites of pairs, to a sector is a convention. In a normal metal $n$ counts electron--hole pairs, and a phonon decays into pairs at the rate $\gamma^{\mathrm{ep}}_{\vect{Q}\nu}$ of \cref{eq:6.4-gamma}; the sectors $(0, m)$ and $(2, m - 1)$ mix within this decay width. In a Dirac semimetal the spectrum is linear and the effective fine-structure constant $\tilde\alpha$ of \cref{sec:7.3} is dimensionless, so the only energy scale near the node is $\omega$ itself. By dimensional analysis the quasiparticle width is then $\omega$ times a function of $\tilde\alpha$, which starts at $\tilde\alpha^2$ because a decay needs two Coulomb vertices. The split is therefore meaningful only for $\tilde\alpha^2 \ll 1$, and near $\tilde\alpha \sim 1$ only totals are quoted. The totals are safe in every case if the response is conserving. The f-sum then holds for the sum of the sectors, while the share of each sector is not fixed. A nonconserving pair-breaking response, such as the bare bubble that violates the f-sum in \cref{sec:7.2}, cannot be paired with any phonon sector.

\begin{workedexample}{A two-level electron on a vibrating bond}\label{we:8-sectors}
We take the bond coordinate $x$ of silicon from \cref{we:5-displaced}, an oscillator of frequency $\omega_0 = \SI{64.1}{meV}$, zero-point length $\ell_0 = \SI{0.048}{\angstrom}$ and mass $M/2$, and attach to it an electron with two levels, a lower one $|\mathrm{lo}\rangle$ and an upper one $|\mathrm{up}\rangle$, split by $\Delta_{\mathrm{tl}} = \SI{1.17}{eV}$, the gap of silicon. With $x = \ell_0(d + d^\dagger)$,
\begin{equation*}
  H = \omega_0\,d^\dagger d + \bigl[\Delta_{\mathrm{tl}} + g_{\mathrm{d}}(d + d^\dagger)\bigr]\,|\mathrm{up}\rangle\langle\mathrm{up}| + g_{\mathrm{n}}(d + d^\dagger)\bigl(|\mathrm{up}\rangle\langle\mathrm{lo}| + |\mathrm{lo}\rangle\langle\mathrm{up}|\bigr) .
\end{equation*}
The diagonal coupling $g_{\mathrm{d}} = -\sqrt2\,\omega_0$ displaces the excited surface as in \cref{we:5-displaced}, with $S_{\mathrm{HR}} = 2$. The off-diagonal coupling $g_{\mathrm{n}} = \SI{20}{meV}$, a model choice, mixes the two levels as the bond moves; it is the only source of nonadiabatic coupling, with the mixing $g_{\mathrm{n}}/\Delta_{\mathrm{tl}} = 0.017$ at the zero-point amplitude. The probes are the electronic transition operator $\sigma_x = |\mathrm{up}\rangle\langle\mathrm{lo}| + |\mathrm{lo}\rangle\langle\mathrm{up}|$ and a kick $\ee^{\ii qx}$ with $q\ell_0 = 1$, that is $q = \SI{41.1}{keV}$, whose recoil energy on the coordinate, $q^2/M$, equals $\omega_0$. We diagonalize $H$ exactly in 160 oscillator states; 90 give the same numbers to the digits quoted.

\emph{Labels.} The two probes reach forty eigenstates with a weight above $10^{-10}$. Each has at least $0.996$ of its norm on a single Born--Oppenheimer product state, $\Psi^{\mathrm{nuc}}_{0,m}(x)\Psi_0(x)$ on the ground surface or $\Psi^{\mathrm{nuc}}_{1,m}(x)\Psi_1(x)$ on the excited one, with $m$ the vibrational level. We assign the ground surface to the electron-following sector, $m = 0$ on the excited surface to the zero-phonon sector and $m \ge 1$ there to the phonon-dressed one. The sector projectors are sums of exact eigenprojectors, they commute with $H$, and the split is exact.

\emph{The electron probe.} The weight of $\sigma_x$, one in total, splits as $1.25\times10^{-3}$, $0.152$ and $0.847$ among the three sectors. The electron-following part is a line at \SI{63.4}{meV}: the bond's ground-state surface, softened by the mixing, carries an electronic dipole $\langle\sigma_x\rangle_0 \simeq -2g_{\mathrm{n}}x/(\ell_0\Delta_{\mathrm{tl}})$ that follows the bond, with the one-phonon weight $(2g_{\mathrm{n}}/\Delta_{\mathrm{tl}})^2 = 1.17\times10^{-3}$ at leading order. The zero-phonon line, at \SI{1.0463}{eV}, has the weight $0.152$, the Franck--Condon factor of the two computed surfaces (0.1519 with the dependence of the transition dipole on $x$). It exceeds the $\ee^{-2} = 0.135$ of \cref{we:5-displaced} because the mixing softens the ground surface to \SI{63.4}{meV} and stiffens the excited one to \SI{65.5}{meV}. The first moment of $\sigma_x$ is fixed by its double commutator at $-\langle(\Delta_{\mathrm{tl}} + g_{\mathrm{d}}x/\ell_0)\sigma_z\rangle = \SI{1.16915}{eV}$, with $\sigma_z = |\mathrm{up}\rangle\langle\mathrm{up}| - |\mathrm{lo}\rangle\langle\mathrm{lo}|$, and the electron-following sector takes \SI{8.32e-5}{eV} of it, a fraction $7.1\times10^{-5}$. At leading order the fraction is $4g_{\mathrm{n}}^2\omega_0/\Delta_{\mathrm{tl}}^3 = 6.4\times10^{-5}$, which is also the drag ratio $\delta M/(M/2)$ of \cref{eq:8.6-fsum-k} for this model, for the reason that made $Z_\kappa m_e/M_\kappa$ appear twice in silicon.

\emph{The kick.} The electrons are left excited with probability $P_{\mathrm{Mig}} = 3.94\times10^{-6}$, of which $0.38\times10^{-6}$ is zero-phonon and $3.56\times10^{-6}$ phonon-dressed. The boost formula \cref{eq:8.6-PMig}, evaluated at $x = 0$ with the vertical gap, gives $(2\omega_0q\ell_0g_{\mathrm{n}}/\Delta_{\mathrm{tl}}^2)^2 = 3.51\times10^{-6}$; the difference between the exact and the boost-formula probabilities, \SI{12}{\percent}, is of the order $S_{\mathrm{HR}}\omega_0/\Delta_{\mathrm{tl}} = 0.11$ of the displacement of the excited surface, which the leading formula ignores. As $g_{\mathrm{n}} \to 0$ the probability falls as $g_{\mathrm{n}}^2$: the kick reaches the electronic sectors through the nonadiabatic coupling alone. The electronic sectors take \SI{4.72}{\micro eV} of the mean recoil energy of \SI{64.1}{meV}, against $P_{\mathrm{Mig}}\,q^2/M + (q^2/M)\,\delta M/(M/2) = 0.25 + \SI{4.10}{\micro eV}$ from \cref{eq:8.6-fsum-k}, and the lattice keeps the rest. The difference of \SI{8}{\percent} between the exact and the predicted energies has the same origin as the \SI{12}{\percent} in $P_{\mathrm{Mig}}$: with the displacement removed, $g_{\mathrm{d}} = 0$, the exact values are $3.52\times10^{-6}$ and \SI{4.34}{\micro eV}, within \SI{0.3}{\percent} of the leading formulas. The energy taken by the electronic sectors is the correction to \cref{eq:5.6-sum} in its smallest form.

\emph{The interference.} The energy-weighted interference of the two probes is $(2.17 - 13.22\,\ii)\times10^{-4}$\,eV in the electron-following sector and exactly the opposite in the two electronic sectors together. Its total vanishes because $\sigma_x$ commutes with $[H, \ee^{\ii qx}]$, the toy form of \cref{eq:8.6-fsum-ek}.

In this model every zero-phonon level is discrete and nothing decays, so the labels are exact and no convention is needed. Couple the excited level to a continuum of phonons, as the band of \cref{we:4-level} couples to its level, and the zero-phonon line acquires a width; the split then follows \cref{eq:8.6-dyson}, with the residue $1.0064$ and the finite-part weight $-0.0064$ of that example.
\end{workedexample}

The decomposition rests on three statements that are not yet theorems. First, two results are derived at the level of perturbation theory in $g$ with distributional limits: the split between the zero-phonon and the phonon-dressed sectors outside the off-shell window, and the convention independence of the totals. A proof needs a limiting-absorption principle for the coupled electron--phonon Hamiltonian at fixed total crystal momentum, on the energies outside the window, uniform in the coupling and in the number of cells. The same estimate would sharpen the gapless bound \cref{eq:8.6-bin} from $(\varGamma\,\Delta E_{\mathrm{c}})^{1/2}$ to the width $\varGamma$. Second, the constants of the superadiabatic construction grow with the number of nuclei, so \cref{eq:8.6-leak} is proved for a finite crystal. The missing step is a superadiabatic generator built from exponentially localized terms, so that the leak per unit volume stays finite in the infinite crystal. Third, for gapless targets it is open whether the dressing $U_{\mathrm{qp}}$ exists beyond perturbation theory, with bounds, since without a gap there is no adiabatic theorem to provide it. The bin bound \cref{eq:8.6-bin} needs this dressing and, with it, a bound on the connected part of $\mathfrak{h}_{\mathrm{off}}$ through a commutator estimate uniform in the number of cells. Until these three statements are proved, the decomposition is exact for a finite insulator and controlled, with the stated error terms, for everything else.
